\section{Two Bar Plane Truss Design problem}

\subsection{Problem Description}
“The Two-Bar Truss problem is composed of two diagonal bars (of aluminium here) connected at a node and fixed at two upper points (Figure \ref{fig:truss_diag}). On the central node, we applied a vertical strength. This force represents the external load that the structure must support. We look for the vertical displacement of the node.”\cite{task3report}

\begin{figure}[H]
    \centering
    \includegraphics[width=0.5\textwidth]{Two Bar Plane Truss.png}
    \caption{Two Bar Plane Truss}
    \label{fig:truss_diag}
\end{figure}

The constitutive equations are \cite{task3report}:
\begin{align}
f_1(x) &= 2 \rho h x_2 \sqrt{1 + x_1^2} \\
f_2(x) &= \frac{P h (1 + x_1^2)^{1.5} \sqrt{1 + x_1^4}}{2 \sqrt{2} E x_1^2 x_2} \\
g_1(x) &= \frac{P(1 + x_1) \sqrt{1 + x_1^2}}{2 \sqrt{2} x_1 x_2} - \sigma_0 \leq 0 \\
g_2(x) &= \frac{P(- x_1 + 1) \sqrt{1 + x_1^2}}{2 \sqrt{2} x_1 x_2} - \sigma_0 \leq 0
\end{align}

\subsection{Method}
The BFGS method needs one function in entry. To obtain the unique formula a weighted sum can be apply. The two function $f_1$ and $f_2$ have the same weight $w_f$ since neither is predominant in solving the problem. To integrate the constraint functions penalties are used. If $g_1$ or $g_2$ are greater than 0, then the penalty equals the value of the function squared. To eliminate solutions that don’t satisfy the constraint functions, a weight $w_c$ is applied. In this way, the value of F is much greater and therefore does not achieve the minimisation objective.
\begin{equation}
F(x) = w_f (f_1(x) + f_2(x)) + w_c (\text{penalty}_{g1} + \text{penalty}_{g2})
\end{equation}

\subsubsection{Exploration of the BFGS parameters}
BFGS optimisation algorithm is governed by a set of main optional parameters set to default. To achieve the best possible optimisation, an exploration of the following parameters is carried out (Table 5 to Table 9). Each series of test is used to create the following in order to benefit from the achieved improvement.

The criteria to choose the best parameter among those tired are the following:
\begin{itemize}
	\item convergence of the algorithm,
	\item low number of iterations,
	\item final gradient norm below a fixed tolerance,
	\item final step size no equal to 0 (from e-6 onwards, this is considered to be 0),
	\item a penalty weight that is not too low to avoid keeping all solutions and not too high to avoid obtaining a solution that is mathematically correct but unfeasible in reality.
\end{itemize}

The maximum number of iterations is a safety limit to prevent infinite loops when convergence is slow or impossible. It should be large enough that algorithm converges naturally via gradient criterion. Each iteration involves computing the gradient, updating the quasi-Hessian approximation, and performing a line search. \cite[Chapter 8]{nocedalwright2006}
The search range is bounded by 50 and 350, and the increment between each test is 100.
\begin{table}[H]
    \centering
    \resizebox{\textwidth}{!}{
    \begin{tabular}{|l|l|l|l|l|l|}
    \hline
    Test & Max iters & Grad norm tol & Step size tol & Fin diff step & Penalty weight \\ \hline
    1 & 50 & 1e-8 & 1e-12 & 1e-6 & 10 000 \\ \hline
    2 & 150 & 1e-8 & 1e-12 & 1e-6 & 10 000 \\ \hline
    3 & 250 & 1e-8 & 1e-12 & 1e-6 & 10 000 \\ \hline
    4 & 350 & 1e-8 & 1e-12 & 1e-6 & 10 000 \\ \hline
    \end{tabular}
    }
    \caption{Exploration of maximum iterations}
    \label{tab:truss_exp_iter}
\end{table}

The gradient norm tolerance is a primary convergence criterion based on the Euclidean norm of the gradient vector. When this norm is inferior to the tolerance, the algorithm terminates as it’s near a stationary point. A small value demands tighter convergence but requires more iterations. \cite[Section 2.1]{nocedalwright2006}
The typical values go from 1e-5 to 1e-8. The search range is bounded by 1e-4 and 1e-8, and the increment between each test is 1e-1.
\begin{table}[H]
    \centering
    \resizebox{\textwidth}{!}{
    \begin{tabular}{|l|l|l|l|l|l|}
    \hline
    Test & Max iters & Grad norm tol & Step size tol & Fin diff step & Penalty weight \\ \hline
    5 & 50 & 1e-4 & 1e-12 & 1e-6 & 10 000 \\ \hline
    6 & 50 & 1e-5 & 1e-12 & 1e-6 & 10 000 \\ \hline
    7 & 50 & 1e-6 & 1e-12 & 1e-6 & 10 000 \\ \hline
    8 & 50 & 1e-7 & 1e-12 & 1e-6 & 10 000 \\ \hline
    \end{tabular}
    }
    \caption{Exploration of gradient norm tolerance}
    \label{tab:truss_exp_grad}
\end{table}

The step size tolerance is a backup convergence criterion when gradient norm tolerance fails. It is based on the change in position between iterations. When parameter updates become below tolerance, the optimisation terminates. It indicates that the algorithm has stalled or converged and prevents unnecessary iterations when making non meaningful progress. \cite[Chapter 2]{nocedalwright2006}
The typical values go from 1e-6 to 1e-9. The search range is bounded by 1e-6 and 1e-12, and the increment between each test is 1e-2.
\begin{table}[H]
    \centering
    \resizebox{\textwidth}{!}{
    \begin{tabular}{|l|l|l|l|l|l|}
    \hline
    Test & Max iters & Grad norm tol & Step size tol & Fin diff step & Penalty weight \\ \hline
    9 & 50 & 1e-5 & 1e-6 & 1e-6 & 10 000 \\ \hline
    10 & 50 & 1e-5 & 1e-8 & 1e-6 & 10 000 \\ \hline
    11 & 50 & 1e-5 & 1e-10 & 1e-6 & 10 000 \\ \hline
    6 & 50 & 1e-5 & 1e-12 & 1e-6 & 10 000 \\ \hline
    \end{tabular}
    }
    \caption{Exploration of step size tolerance}
    \label{tab:truss_exp_step}
\end{table}

The finite difference step is a gradient approximation accuracy. If it’s too large, it creates a truncation error. If it’s too small, it creates a roundoff error due to floating-point arithmetic. \cite[Section 7.1]{nocedalwright2006}
The typical values go from 1e-5 to 1e-8. The search range is bounded by 1e-4 and 1e-8, and the increment between each test is 1e-1.
\begin{table}[H]
    \centering
    \resizebox{\textwidth}{!}{
    \begin{tabular}{|l|l|l|l|l|l|}
    \hline
    Test & Max iters & Grad norm tol & Step size tol & Fin diff step & Penalty weight \\ \hline
    12 & 50 & 1e-5 & 1e-12 & 1e-4 & 10 000 \\ \hline
    13 & 50 & 1e-5 & 1e-12 & 1e-5 & 10 000 \\ \hline
    6 & 50 & 1e-5 & 1e-12 & 1e-6 & 10 000 \\ \hline
    14 & 50 & 1e-5 & 1e-12 & 1e-7 & 10 000 \\ \hline
    15 & 50 & 1e-5 & 1e-12 & 1e-8 & 10 000 \\ \hline
    \end{tabular}
    }
    \caption{Exploration of finite difference step}
    \label{tab:truss_exp_fd}
\end{table}

The penalty weight controls the trade-oof between optimising the original objective and enforcing constraint satisfaction. A low penalty weight allows the algorithm to find better objective values but with significant violations. A high penalty weight strongly enforced the constraint satisfaction, but the problem becomes ill-conditioned and BFGS convergence may suffer.
The typical values can be between 1 to 100. It should be adapted to the problem. The search range is bounded by 100 and 100 000, and the increment between each test is x10.
\begin{table}[H]
    \centering
    \resizebox{\textwidth}{!}{
    \begin{tabular}{|l|l|l|l|l|l|}
    \hline
    Test & Max iters & Grad norm tol & Step size tol & Fin diff step & Penalty weight \\ \hline
    16 & 50 & 1e-5 & 1e-12 & 1e-6 & 100 \\ \hline
    17 & 50 & 1e-5 & 1e-12 & 1e-6 & 1000 \\ \hline
    6 & 50 & 1e-5 & 1e-12 & 1e-6 & 10 000 \\ \hline
    18 & 50 & 1e-5 & 1e-12 & 1e-6 & 100 000 \\ \hline
    19 & 50 & 1e-5 & 1e-12 & 1e-6 & 1 000 000 \\ \hline
    \end{tabular}
    }
    \caption{Exploration of penalty weights}
    \label{tab:truss_exp_penalty}
\end{table}

\subsubsection{Exploration of noise effects}
The objective of this study is to evaluate the impact of noise on the convergence of BFGS, applied to the Two-Bar Truss problem, formulated via a weighted objective function with constraint penalties.

The results show that noise strongly influences the behaviour of the solver, not only on the quality of the solution, but above all on the reliability of the convergence criteria.

\begin{table}[H]
    \centering
    \resizebox{\textwidth}{!}{
    \begin{tabular}{|l|l|}
    \hline
    Noise Level & Behaviour \\ \hline
    Low Noise ($\sigma \leq 10^{-6}$) & Reliable convergence, optimal solutions \\ \hline
    Moderate noise ($10^{-5} \leq \sigma \leq 10^{-3}$) & Interesting compromise: improved feasibility, faster convergence \\ \hline
    Loud noise ($\sigma \geq 10^{-2}$) & Apparent convergence but degraded solutions \\ \hline
    \end{tabular}
    }
    \caption{Noise Classification}
    \label{tab:truss_noise_class}
\end{table}

To begin we can start with the noise application:
\begin{equation}
f_{noisy}(x) = f_{clean}(x) + \sigma \mathcal{N}(0,1)
\end{equation}
\begin{itemize}
    \item $f_{clean}(x)$: deterministic objective function (weighted sum of objectives and penalty terms)
    \item $\sigma$: noise level (noise amplitude)
    \item $\mathcal{N}(0,1)$: standard Gaussian noise (zero mean, unit variance)
    \item The noise is independent at each function evaluation
\end{itemize}

In this work, noise is added directly to the objective function evaluation. This noise can be seen as a measurement error, like a sensor defect or inaccuracy in evaluating system performance. In practice, the values used for optimisation are not always perfectly accurate, due to noisy measurements or numerical errors. Adding this noise therefore makes it possible to approximate more realistic conditions and study how the BFGS algorithm reacts to these uncertainties.

To evaluate the quality of the solutions obtained, several acceptance criteria are used. A solution is first considered acceptable if the BFGS algorithm terminates correctly according to its internal criteria. Next, the solution must comply with the mechanical constraints of the problem, i.e. the tensile and compressive constraints must be satisfied. Finally, the quality of the solution is evaluated using the noise-free objective function to verify that the solution found remains effective despite the disturbances introduced during optimisation.

The first acceptance criterion is based on the internal criteria of the BFGS algorithm. A solution is considered to have converged when the algorithm stops according to its stopping conditions, such as a very small variation in the solution between two iterations, a sufficiently small gradient norm, or the maximum number of iterations being reached. However, in the presence of noise, this convergence is primarily numerical and does not necessarily guarantee that the solution is optimal from a physical point of view.

The second criterion concerns compliance with the mechanical constraints of the problem. A solution is said to be feasible only if the tensile and compressive stresses are satisfied, i.e. if the associated stresses remain less than or equal to zero. These stresses are evaluated without noise, based on the physical expressions of the problem, to ensure that the structure complies with the material's strength limits.

Finally, the quality of the solution is evaluated using the deterministic objective function, without noise. This re-evaluation makes it possible to compare the actual performance of the solutions found and to verify that they remain effective despite the perturbations introduced during optimisation. This criterion is essential for distinguishing a truly interesting solution from one that would appear to be optimal solely because of noise.

The results show that noise strongly influences the behaviour of the solver.
\begin{figure}[H]
    \centering
    \includegraphics[width=0.8\textwidth]{Noise effect on the Two Truss Problem.png}
    \caption{Noise effect on the Two Bar Truss problem}
    \label{fig:truss_noise}
\end{figure}

\subsection{Results}

\subsubsection{Exploration of the BFGS parameters}
NB: The results of each simulation are available in the Appendix A.A.

The first parameter exploration is the one about maximum iterations. The tests 1 to 4 (see section 5.2.1) didn’t converge (see Appendix A.A), have the same number of iterations and produce the same results. Considering this information and that the number of iterations is 30 (Figure 15 a) the maximum iterations parameter chosen is 50. It generates non consistent results (Figure 15 b and c).

\begin{figure}[H]
    \centering
    \begin{subfigure}[b]{0.48\textwidth}
        \includegraphics[width=\textwidth]{BFGS Convergence objective function Max iter.png}
        \caption{Objective Function}
        \label{fig:truss_exp_obj}
    \end{subfigure}
    \hfill
    \begin{subfigure}[b]{0.48\textwidth}
        \includegraphics[width=\textwidth]{BFGS Convergence gradient Max iter.png}
        \caption{Gradient Norm}
        \label{fig:truss_exp_grad}
    \end{subfigure}
    \hfill
    \begin{subfigure}[b]{0.48\textwidth}
        \includegraphics[width=\textwidth]{BFGS Step sizes Max iter.png}
        \caption{Step Sizes}
        \label{fig:truss_exp_step}
    \end{subfigure}
    \caption{BFGS Parameter Exploration Results}
    \label{fig:truss_exp_all}
\end{figure}

The exploration of the following parameters gradient norm tolerance, step size tolerance, finite difference step, and penalty weight conducted to confirm the first choice of values.
Tests 5 and 6 have same results and converged on the contrary of the previous simulation. The choice is test 6 because more conservative so safer for other problems.
Testing step size tolerance showed no performance difference. All configurations converged identically in 24 iterations via the gradient criterion. The step tolerance did not trigger because the algorithm maintained healthy steps throughout, demonstrating that for well-behaved problems, gradient tolerance is the primary convergence criterion.
Testing finite difference step from 1e-4 to 1e-8 revealed that 1e-4 is too big, requiring 25\% more iterations (30 vs 24) and yielding poor gradient accuracy (1.2e-03 vs 1.2e-06). Values from 1e-5 to 1e-8 performed identically. The optimal value 1e-6 balances truncation error and roundoff error.
The analyse of the influence of the penalty weight gives the following indications that conducted to selec a penalty of 10 000.
\begin{itemize}
	\item 100 and 1000: too small, constrained ignored, physically unrealistic
	\item 10 000: very good default
	\item 100 000: stricter but still reasonable
	\item 1 000 000: too large, weighted sum irrelevant, only feasibility matters, distorts search
\end{itemize}

Taking into account previous simulations, the set of parameters to obtain the final result (Figure \ref{fig:truss_final_obj} to \ref{fig:truss_final_step}) is:
\begin{itemize}
	\item Algorithm parameters:
		\begin{itemize}
			\item Maximum number of iterations: 50
			\item Gradient norm tolerance: 1e-5
			\item Step size tolerance: 1e-12
			\item Finite difference step: 1e-6
		\end{itemize}
	\item Problem parameters:
		\begin{itemize}
			\item Penalty weights: 10 000
		\end{itemize}
\end{itemize}

\begin{figure}[H]
    \centering
    \includegraphics[width=0.8\textwidth]{BFGS Convergence objective function Best settings.png}
    \caption{BFGS convergence for the objective function}
    \label{fig:truss_final_obj}
\end{figure}
The penalised objective function F(x) decreases over 24 iterations. The rapid initial descent (iterations 0 to 5) indicates the algorithm quickly moves away from the infeasible initial guess toward the feasible region. The convergence slows down around 15 iterations and after as the algorithm approaches the optimal solution.

\begin{figure}[H]
    \centering
    \includegraphics[width=0.8\textwidth]{BFGS Convergence gradient Best settings.png}
    \caption{BFGS convergence for the gradient norm}
    \label{fig:truss_final_grad}
\end{figure}
The gradient norm decreases approximately from 1e2 to 1e-6 satisfying the gradient norm tolerance of 1e-5. The behaviour is non-monotonic, before iteration 19, reflects the algorithm’s adjustment as it balances the two objectives $f_1$ and $f_2$ and constraint penalties. The final gradient norm of 1.20e-6 (rounded to two decimal places) confirms convergence to a stationary point.

\begin{figure}[H]
    \centering
    \includegraphics[width=0.8\textwidth]{BFGS Step sizes Best settings.png}
    \caption{BFGS step sizes}
    \label{fig:truss_final_step}
\end{figure}
Step sizes remain consistently at 1.0 (full step) throughout the optimisation after iteration 15, indicating that the BFGS approximation and line search are performing will. The algorithm doesn’t need to reduce the step lengths to ensure the improvement in optimisation. This suggests that the problem is well-conditioned. Moreover, the consistent unit step length is characteristic of well-behaved quasi-Newton methods.

With this simulation the two bar plane truss characteristics are optimised and give the following result (Table 11).
\begin{table}[H]
    \centering
    \begin{tabular}{|l|l|l|l|l|l|}
    \hline
    x1 [norm] & Position & x2 [norm] & Area & Weight & Displacement \\ \hline
    0.66 & 1.67 m & 0.53 & 3.44 $m^2$ & 16.34 kg & 8.67e-3 m \\ \hline
    \end{tabular}
    \caption{Optimised two bar plane truss (rounded to two decimal places)}
    \label{tab:truss_optimal}
\end{table}

\subsubsection{Exploration of noise effects}

We can saw in the Figure \ref{fig:truss_noise} Noise effect on the Two Bar Truss problem, the BFGS Convergence vs Noise figure shows the evolution of the algorithm's convergence rate as a function of the noise level applied to the objective function. We observe that the convergence rate remains above 50\% overall (the convergence rate is the proportion of simulations for which the BFGS algorithm terminates by declaring convergence), indicating that the algorithm still manages to terminate according to its internal criteria. However, this evolution is not monotonic: some noise levels seem to slightly improve convergence, while others degrade it. This highlights the fact that noise can sometimes facilitate the algorithm's progress but also disrupt its behaviour. These results show that the numerical convergence criterion alone is not sufficient to assess the actual quality of a solution when noise is present. We can saw that in the follow figure.

\begin{figure}[H]
    \centering
    \includegraphics[width=0.8\textwidth]{The solution quality.png}
    \caption{The solution quality}
    \label{fig:truss_sol_qual}
\end{figure}

The figure Solution Quality vs Noise illustrates the impact of noise on the quality of the solutions obtained by the BFGS algorithm. Although the algorithm often converges according to its numerical stopping criteria, the value of the deterministic objective function evaluated after optimisation shows an increasing dispersion as the noise level rises. This indicates that some solutions, while numerically converged, remain far from optimal when assessed without noise. These results confirm that numerical convergence alone is not sufficient to properly assess solution quality in the presence of noise.

Beyond the dispersion observed in the objective function values, it is also essential to verify whether the solutions obtained remain physically acceptable. Indeed, a solution that appears satisfactory in terms of objective value may still violate the mechanical constraints of the problem. For this reason, the analysis is extended by examining the feasibility of the solutions with respect to the tension and compression constraints, as illustrated in the following figure.
\begin{figure}[H]
    \centering
    \includegraphics[width=0.8\textwidth]{Feasability of the solution.png}
    \caption{Feasibility of the solutions}
    \label{fig:truss_feas}
\end{figure}

This graph (Figure 17) shows the evolution of the feasible solution rate as a function of the level of noise added to the objective function. A solution is considered feasible when the tension and compression constraints are met. We observe that, for very low noise levels, a large proportion of the solutions obtained do not meet the constraints, even though the algorithm converges numerically. As the noise increases, the rate of feasible solutions improves significantly and exceeds the 80\% threshold. This shows that numerical convergence alone is not sufficient to guarantee a physically acceptable solution, and that it is necessary to explicitly verify compliance with the constraints, particularly in the presence of noise.

\subsection{Discussion}
\subsubsection{Parameter exploration results discussion}
The exploration of BFGS parameters revealed sensitivities and optimal ranges for the Two Bar Plane Truss problem. 

Tests 1 to 4 demonstrated that maximum iterations above 50 provide no benefit (all converged between 24 and 30 iterations with identical results). This confirms that convergence is governed by gradient criteria rather than iteration limits. Additionally, the value 50 balances safety and computational cost.
For gradient norm tolerance, 1e-5 (test 6) was selected over 1e-4 (test 5) for conservatism despite identical behaviour. Tighter tolerances (tests 7 and 8) prevented convergence. The final gradient norm of 1.20e-6 confirms the algorithm achieved precision beyond the specified tolerance.
Simulations 9 to 11 and 6 revealed that step size tolerance has negligible impact on algorithm performance. Indeed, all configurations converged in 24 iterations and had a final step size of 1.0. This indicates the problem is well-conditioned and gradient-based convergence dominates. Therefore, it can be concluded that the step tolerance is a backup criterion.
Test 12 demonstrated that a low finite difference step creates excessively large perturbations that degrading the performance. On the contrary, tests 13 to 15 performed nearly identically. The selected value of 1e-6 optimally balances truncation error and roundoff error.
The exploration of the penalty weight (tests 16 to 19) revealed a critical sensitivity. Insufficient penalties allowed constraint violation leading to physical infeasibility (test 16), premature convergence (test 17), non-convergence (test 18) and distorting the objective function landscape (test 19). The selected penalty weight of 10 000 (test 6) ensures physical feasibility while maintaining numerical stability.

Quasi-Newton BFGS method is sensitive to its parameterisation. Therefore, it is important to choose it carefully by realising several simulations to obtain the best optimisation.

These findings align with theoretical predictions for quasi-Newton methods on well-conditioned problems \cite[Chapter 8]{nocedalwright2006}. The gradient-dominated convergence behaviour demonstrates the algorithm’s efficiency when applied to smooth and differentiable objectives. The sensitivity to penalty weight reflects conditioning issues discussed by Shanno \cite{shanno1970}, where ill-conditioned problems can slow or prevent convergence.

\subsubsection{Noise exploration results discussion}
We can saw in the Figure \ref{fig:truss_noise} Noise effect on the Two Bar Truss problem, the BFGS Convergence vs Noise figure shows the evolution of the algorithm's convergence rate as a function of the noise level applied to the objective function. We observe that the convergence rate remains above 50\% overall (the convergence rate is the proportion of simulations for which the BFGS algorithm terminates by declaring convergence), indicating that the algorithm still manages to terminate according to its internal criteria. However, this evolution is not monotonic: some noise levels seem to slightly improve convergence, while others degrade it. This highlights the fact that noise can sometimes facilitate the algorithm's progress but also disrupt its behaviour. These results show that the numerical convergence criterion alone is not sufficient to assess the actual quality of a solution when noise is present. We can saw that in the follow figure.

The systematic exploration of noise levels highlighted the sensitivity of the BFGS algorithm to perturbations in the objective function. While BFGS performs efficiently in a deterministic setting, the presence of noise degrades gradient reliability and leads to increased dispersion in the final solutions.

This table presents the final solutions obtained for different noise levels. It is important to note that only valid solutions are reported, i.e. solutions that have converged and comply with the tensed and compressive constraints.

\begin{table}[H]
    \centering
    \resizebox{\textwidth}{!}{
    \begin{tabular}{|l|l|l|l|l|}
    \hline
    Noise $\sigma$ & x1 (normalised) & x2 (normalised) & Weight [kg] & Displacement [m] \\ \hline
    0.00E+00 & 0.6573 & 0.5328 & 16.3670 & 8.6692e-03 \\ \hline
    1.00E-06 & 0.6572 & 0.5328 & 16.3672 & 8.6693e-03 \\ \hline
    1.00E-05 & 0.6581 & 0.5326 & 16.3672 & 8.6680e-03 \\ \hline
    1.00E-04 & 0.6327 & 0.5394 & 16.3871 & 8.6983e-03 \\ \hline
    1.00E-03 & 0.6094 & 0.5468 & 16.4383 & 8.7188e-03 \\ \hline
    1.00E-02 & 0.5585 & 0.5646 & 16.6031 & 8.7782e-03 \\ \hline
    1.00E-01 & 0.5692 & 0.5608 & 16.5659 & 8.7616e-03 \\ \hline
    1.00E+00 & 0.2647 & 1.4213 & 37.7449 & 5.6638e-03 \\ \hline
    \end{tabular}
    }
    \caption{Noise exploration results}
    \label{tab:truss_noise_exp}
\end{table}

For low noise levels, the design variables and physical performance remain very close to the reference solution without noise. As noise increases, the solutions become progressively more dispersed and conservative, with an increase in weight and displacement. At high noise levels, although the solutions remain mechanically acceptable, their quality deteriorates significantly, highlighting the impact of noise on the effectiveness of the BFGS algorithm.

\subsubsection{Work discussion}
These experimentations allow us to understand better the BFGS parameter sensitivity and constraint handling via penalty methods. We also learn about the BFGS noisy comportment and the unexpected effects on the results.

The systematic parameter exploration methodology effectively identified optimal values and revealed sensitivities. The BFGS’s gradient-based approach required careful problem formulation that is suitable for quasi-Newton optimisation. \cite{fletcher1970, shanno1970, nocedalwright2006}
Future work should include a sensitivity analysis about how different starting points affect convergence speed or risk of local minima. Moreover, the penalty method has limitations. An adaptative method could provide more robust constraint handling. Finally, the algorithm verifies first-order condition (gradient norm) but doesn’t check second-order sufficient conditions (Hessian matrix). It would be necessary to implement it to confirm the solution is a true local minimum rather than a saddle point.
The systematic exploration of noise levels highlighted the sensitivity of the BFGS algorithm to perturbations in the objective function. While BFGS performs efficiently in a deterministic setting, the presence of noise degrades gradient reliability and leads to increased dispersion in the final solutions. This confirms that gradient-based quasi-Newton methods require careful validation when applied to noisy optimisation problems.

Future work could extend this study by investigating the impact of input data truncation on the optimisation process. In practical applications, measurements or simulations are often limited in precision, which could further affect the behaviour of gradient-based methods. In addition, since the Two-Bar Truss problem allows an explicit analytical formulation, it would be possible to implement a full Newton method using the true Hessian matrix. This would enable a direct comparison with the quasi-Newton BFGS approach and provide further insight into the influence of second-order information on convergence robustness and solution quality.

\subsection{Comparison with the task 3}

\begin{table}[H]
	\centering
	\resizebox{\textwidth}{!}{
		\begin{tabular}{|l|l|l|l|l|}
			\hline
			Algorithm & x1 (normalised) & x2 (normalised) & Weight [kg] & Displacement [m] \\ \hline
			PSO & 1.2560 & 1.1667 & 48.0982 & 0.0029 \\ \hline
			SA & 1.0429 & 1.0043 & 37.2536 & 0.0038 \\ \hline
			BFGS & 0.6573 & 0.5328 & 16.3670 & 0.0087 \\ \hline
		\end{tabular}
	}
	\caption{Results (rounded to four decimal) comparison \cite{task3report} (see section 3.2)}
\end{table}

The BFGS optimization achieved significantly superior results compared to PSO and SA algorithms \cite{task3report}. BFGS produced a structure weighing 16.37 kg, representing a 66\% reduction compared to PSO (48.10 kg) and 56\% reduction compared to SA (37.25 kg). This demonstrates BFGS's ability to exploit local gradient information for precise convergence, as predicted by quasi-Newton theory \cite[Chapter 8]{nocedalwright2006}. The gradient-based approach proves superior for this deterministic, smooth, constrained optimization problem, achieving convergence in only 24 iterations compared to the population-based evaluations required by metaheuristics.